% GENERATED FILE. Edit canonical lessons, then run npm run generate:books.
% canonical-source-hash: fnv1a64:42d5412a0984511e
% canonical-lessons: MR-C196-phirayla-jane, MR-C196-ghadi-karne, MR-C196-tayari-karne, MR-C196-olakh-karun-dene, MR-C196-usne-dene

\chapter{To go for a walk, To fold, To prepare, To introduce, To lend}
\label{ch:mr-to-go-for-a-walk-to-fold-to-prepare-to-introduce-to-lend}
\hlchaptermodality{196}

\begin{chapteropening}
\textbf{By the end of this chapter:} \emph{I can say to go for a walk, to fold, to prepare, to introduce and to lend.}

Complete the last lesson of chapter 196: I can say to go for a walk, to fold, to prepare, to introduce and to lend.
\end{chapteropening}

\section[\texorpdfstring{\mr{फिरायला} \mr{जाणे} (phirāylā jāṇe)}{फिरायला जाणे (phirāylā jāṇe)}]{\mr{फिरायला} \mr{जाणे} (phirāylā jāṇe) — to go for a walk}

\label{lesson:MR-C196-phirayla-jane}



\begin{quote}
Before the new one: say the Marathi for to knock, then the Marathi for to compare.
\end{quote}

\subsection*{You'll want to know: \mr{फिरायला} \mr{जाणे}}
\textbf{\mr{फिरायला} \mr{जाणे}} — \emph{phirāylā jāṇe} — \textquotedblleft{}to go for a walk\textquotedblright{}.

\begin{grammarlens}[title={What you've built}]
One more word for this chapter.
\end{grammarlens}

\subsection*{Your turn}
\begin{itemize}
\raggedright
  \item \emph{Say it:} \emph{phirāylā jāṇe}
  \item \emph{Say it:} \emph{phirāylā jāṇe}, once more
  \item \emph{Say it:} say \emph{phirāylā jāṇe}
  \item \emph{Recall:} say the Marathi for to knock, then the Marathi for to compare, then say \emph{phirāylā jāṇe} again
\end{itemize}

\begin{tcolorbox}[breakable,colback=teal!4,colframe=teal!35!black,title={Before you move on}]
What does \mr{फिरायला} \mr{जाणे} mean? (\textquotedblleft{}To go for a walk\textquotedblright{}.) Say it once more.
\end{tcolorbox}
\section[\texorpdfstring{\mr{घडी} \mr{करणे} (ghaḍī karṇe)}{घडी करणे (ghaḍī karṇe)}]{\mr{घडी} \mr{करणे} (ghaḍī karṇe) — to fold}

\label{lesson:MR-C196-ghadi-karne}



\begin{quote}
Before the new one: say the Marathi for to compare, then the Marathi for to go for a walk.
\end{quote}

\subsection*{You'll want to know: \mr{घडी} \mr{करणे}}
\textbf{\mr{घडी} \mr{करणे}} — \emph{ghaḍī karṇe} — \textquotedblleft{}to fold\textquotedblright{}.

\begin{grammarlens}[title={What you've built}]
One more word for this chapter.
\end{grammarlens}

\subsection*{Your turn}
\begin{itemize}
\raggedright
  \item \emph{Say it:} \emph{ghaḍī karṇe}
  \item \emph{Say it:} \emph{ghaḍī karṇe}, once more
  \item \emph{Say it:} say \emph{ghaḍī karṇe}
  \item \emph{Recall:} say the Marathi for to compare, then the Marathi for to go for a walk, then say \emph{ghaḍī karṇe} again
\end{itemize}

\begin{tcolorbox}[breakable,colback=teal!4,colframe=teal!35!black,title={Before you move on}]
What does \mr{घडी} \mr{करणे} mean? (\textquotedblleft{}To fold\textquotedblright{}.) Say it once more.
\end{tcolorbox}
\section[\texorpdfstring{\mr{तयारी} \mr{करणे} (tayārī karṇe)}{तयारी करणे (tayārī karṇe)}]{\mr{तयारी} \mr{करणे} (tayārī karṇe) — to prepare}

\label{lesson:MR-C196-tayari-karne}



\begin{quote}
Before the new one: say the Marathi for to go for a walk, then the Marathi for to fold.
\end{quote}

\subsection*{You'll want to know: \mr{तयारी} \mr{करणे}}
\textbf{\mr{तयारी} \mr{करणे}} — \emph{tayārī karṇe} — \textquotedblleft{}to prepare\textquotedblright{}.

\begin{grammarlens}[title={What you've built}]
One more word for this chapter.
\end{grammarlens}

\subsection*{Your turn}
\begin{itemize}
\raggedright
  \item \emph{Say it:} \emph{tayārī karṇe}
  \item \emph{Say it:} \emph{tayārī karṇe}, once more
  \item \emph{Say it:} say \emph{tayārī karṇe}
  \item \emph{Recall:} say the Marathi for to go for a walk, then the Marathi for to fold, then say \emph{tayārī karṇe} again
\end{itemize}

\begin{tcolorbox}[breakable,colback=teal!4,colframe=teal!35!black,title={Before you move on}]
What does \mr{तयारी} \mr{करणे} mean? (\textquotedblleft{}To prepare\textquotedblright{}.) Say it once more.
\end{tcolorbox}
\section[\texorpdfstring{\mr{ओळख} \mr{करून} \mr{देणे} (oḷakh karūn deṇe)}{ओळख करून देणे (oḷakh karūn deṇe)}]{\mr{ओळख} \mr{करून} \mr{देणे} (oḷakh karūn deṇe) — to introduce}

\label{lesson:MR-C196-olakh-karun-dene}



\begin{quote}
Before the new one: say the Marathi for to fold, then the Marathi for to prepare.
\end{quote}

\subsection*{You'll want to know: \mr{ओळख} \mr{करून} \mr{देणे}}
\textbf{\mr{ओळख} \mr{करून} \mr{देणे}} — \emph{oḷakh karūn deṇe} — \textquotedblleft{}to introduce\textquotedblright{}.

\begin{grammarlens}[title={What you've built}]
One more word for this chapter.
\end{grammarlens}

\subsection*{Your turn}
\begin{itemize}
\raggedright
  \item \emph{Say it:} \emph{oḷakh karūn deṇe}
  \item \emph{Say it:} \emph{oḷakh karūn deṇe}, once more
  \item \emph{Say it:} say \emph{oḷakh karūn deṇe}
  \item \emph{Recall:} say the Marathi for to fold, then the Marathi for to prepare, then say \emph{oḷakh karūn deṇe} again
\end{itemize}

\begin{tcolorbox}[breakable,colback=teal!4,colframe=teal!35!black,title={Before you move on}]
What does \mr{ओळख} \mr{करून} \mr{देणे} mean? (\textquotedblleft{}To introduce\textquotedblright{}.) Say it once more.
\end{tcolorbox}
\section[\texorpdfstring{\mr{उसने} \mr{देणे} (usne deṇe)}{उसने देणे (usne deṇe)}]{\mr{उसने} \mr{देणे} (usne deṇe) — to lend}

\label{lesson:MR-C196-usne-dene}



\begin{quote}
Before the new one: say the Marathi for to prepare, then the Marathi for to introduce.
\end{quote}

\subsection*{You'll want to know: \mr{उसने} \mr{देणे}}
\textbf{\mr{उसने} \mr{देणे}} — \emph{usne deṇe} — \textquotedblleft{}to lend\textquotedblright{}.

\begin{grammarlens}[title={What you've built}]
One more word for this chapter.
\end{grammarlens}

\subsection*{Your turn}
\begin{itemize}
\raggedright
  \item \emph{Say it:} \emph{usne deṇe}
  \item \emph{Say it:} \emph{usne deṇe}, once more
  \item \emph{Say it:} say \emph{usne deṇe}
  \item \emph{Recall:} say the Marathi for to prepare, then the Marathi for to introduce, then say \emph{usne deṇe} again
\end{itemize}

\begin{tcolorbox}[breakable,colback=teal!4,colframe=teal!35!black,title={Before you move on}]
What does \mr{उसने} \mr{देणे} mean? (\textquotedblleft{}To lend\textquotedblright{}.) Say it once more.
\end{tcolorbox}
